\documentclass[11pt]{article}

\usepackage[margin=1in]{geometry}
\usepackage{amsmath,amssymb,amsthm,mathtools}
\usepackage{enumitem}
\usepackage{booktabs}
\usepackage{tabularx,longtable,array}
\usepackage{microtype}
\usepackage[hidelinks]{hyperref}
\usepackage{xurl}
\usepackage[nameinlink,noabbrev]{cleveref}


\newcommand{\MaxL}{\mathcal{L}}

\newcommand{\TV}{\mathrm{TV}}

\newcommand{\PP}{\mathbb{P}}

\title{Worked Example: Instantiating a Tiered Reader-Privacy Receipt\\
\large Receipt line items, user discovery, and evaluation interlocks}
\author{}
% NOTE: The shipped worked-example artifacts (receipt, TW/ENF/state registries, validation report)
% correspond to note_version=0.59 and release_id=worked-example-draft-59. This paper is an
% explanatory wrapper around that bundle and should not drift from the artifact ids.
\date{Unified archive note v0.69 (2026-03-05; archive v162; worked-example-draft-59)}


\begin{document}
\maketitle

\begin{abstract}
The synthesis layer provides clean interfaces for receipt line items, user-verifiability, threat windows, exposure normal forms, and auditor replay hooks.
This note is a single end-to-end instantiation whose job is \emph{coherence}, not novelty:
a reader-privacy claim with a relayed delegated-routing primary path and a budgeted structured-overlay fallback.
It is an ``integration test'' for the publication surface: a referee should be able to understand \emph{what is being claimed}, a user should be able to \emph{diff the claim across releases}, and an auditor should be able to \emph{replay} the declared computations after requesting the private artifact bundle.
\end{abstract}

\paragraph{How to read this note.}
Skip everything except:
(i) \Cref{sec:scenario} for the declared threat window and exposure surfaces,
(ii) \Cref{sec:lineitems} for the five representative receipt line items and their replay hooks,
and (iii) \Cref{sec:numerical} for the tiny arithmetic slice that downstream tools consume.
All other sections are there only to show the interfaces compose without ``semantic drift''.

\paragraph{Owned background (avoid redundancy).}
Schema-level receipt tuples and JSONPath meanings are owned by Synthesis~12.\cite{series:receiptitems}
Exposure normal forms (ENF IDs) and observation tiers are owned by Synthesis~3 and Synthesis~15.\cite{series:traceifaces,series:tieredobs}
Threat windows (TW IDs) are owned by Synthesis~4.\cite{series:threatwindows}
Primary-path separation guards are owned by Synthesis~22--23.\cite{series:primarysepmanifest,series:probsep}
Replay-plan minimality and plan/state equivalence are owned by Synthesis~18--20.\cite{series:planstateequiv,series:auditorreplay}
User discovery and minimal interfaces are owned by Synthesis~13.\cite{series:userminiface}
Endpoint conversions are owned by Synthesis~5 (and odds inflation by Anonymity~A).\cite{series:endpointbridge,series:oddsinflation}
Artifact availability and disclosure are governed by ABOM/OINL and the series artifact policy.\cite{series:abomoinl,series:artifactpolicy}

\section{Scenario and declared interfaces}
\label{sec:scenario}
We instantiate one deployment slice with two paths:
\begin{enumerate}[leftmargin=*]
\item \textbf{Primary: relayed delegated routing.} A relay separates client-side metadata from destination-side metadata (a primary separation claim guarded by a published failure probability $\rho$). Deployed control-plane surfaces for delegated routing (AutoConf, overrides, path scopes) are pinned by the state declaration per Synthesis~30.\cite{series:primarysepmanifest,series:deployedsurfaces}
\item \textbf{Fallback: structured overlay lookup.} When primary fails (or is bypassed by policy), a structured overlay lookup is executed and may expose a Tier~2 committee-contact trace (CCT) together with coarse timing and congestion witnesses.\cite{series:traceifaces,series:tieredobs}
\end{enumerate}
The deployer publishes a small \emph{public bundle} (receipt + digest-bound companion objects) and retains a \emph{private bundle} (raw logs, scripts, traces, and any privileged measurements).
Users do \emph{not} need the private bundle to understand or compare the claim; auditors can request it to replay the declared computations.\cite{series:auditorreplay}

\paragraph{Primary-path specificity without bloat.}
The public \texttt{primary\_surface\_manifest\_id} is intentionally a \emph{minimal guard object}: it commits to the separation class and the clauses that would invalidate the non-collusion assumption.\cite{series:primarysepmanifest}
Protocol- and endpoint-level details (e.g., the delegated-routing API version, endpoint set, service discovery/AutoConf mode, and any caching knobs) live in the private bundle as a config snapshot referenced by the manifest's evidence hooks; this keeps the public receipt diffable without forcing a full configuration dump into the paper.
For the deployed-surface taxonomy and a ``what-to-pin'' checklist (projection vs state vs separation), see Synthesis~30.\cite{series:deployedsurfaces}

\paragraph{Where the worked bundle lives.}
In the archive this example is represented by the receipt file \texttt{artifacts/example\_receipt.json} together with digest-bound companion objects.
The full file list is intentionally not reproduced here; it is discoverable from \texttt{support\_manifest.json} and \texttt{example\_artifact\_inventory.json} in the bundle.

\begin{table}[t]
\centering
\small
\begin{tabularx}{\linewidth}{@{}lX@{}}
\toprule
Public object (digest-bound) & What it fixes (why a user can diff it)\\
\midrule
\texttt{receipt\_id} & The line items: declared interfaces, budgets, and replay hooks (Synthesis~12).\\
\texttt{tw\_id} & Threat window declaration: secret, repetition unit/window, adversary model (Synthesis~4).\\
\texttt{enf\_registry\_id} & Exposure normal form labels: which transcript projection is in scope (Synthesis~3).\\
\texttt{state\_decl\_id} & The small public state surface that conditions the claim (state series + Synthesis~18).\\
\texttt{primary\_surface\_manifest\_id} & Guard object for non-collusion/separation assumptions (Synthesis~22).\\
\texttt{replay\_plans\_id} & Auditor replay recipe: how to recompute each line item from private artifacts (Synthesis~20).\\
\bottomrule
\end{tabularx}
\caption{Minimal public identifiers for the worked example. The point is not the particular files but the \emph{digest discipline}: a name like \texttt{enf.fallback.cct.v1} is not trusted unless its digest matches.}
\label{tab:min-public-ids}
\end{table}

\section{Receipt line items}
\label{sec:lineitems}
This note does not re-specify the receipt schema; it \emph{instantiates} it.
The canonical tuple schema lives in Synthesis~12, and deployed control-plane surfaces referenced by any line item are owned by Synthesis~30.\cite{series:receiptitems,series:deployedsurfaces}
Here we keep only the slice needed to make the scenario mechanically replayable.

\paragraph{Minimal slice used here.}
\begin{enumerate}[leftmargin=*]
\item \textbf{Interface IDs:} \texttt{tw\_id} (threat window template) and \texttt{exposure\_nf\_id} (declared observable projection / ENF)---owned by Synthesis~4 and Synthesis~3.\cite{series:threatwindows,series:traceifaces}
\item \textbf{State binding:} \texttt{state\_decl\_id} pins the state surfaces this bundle admits (e.g., the public refresh timer and bucket-level staleness) and therefore what ``replay'' means.\cite{series:planstateequiv}
\item \textbf{Primary separation:} \texttt{primary\_path\_declaration} publishes the separation class and its failure guard $(0,\rho)$.\cite{series:primarysepmanifest,series:probsep}
\item \textbf{Fallback exposure:} \texttt{fallback\_contact\_surface} and \texttt{fallback\_tiered\_vector} publish attenuated Tier~2 summaries under an explicit activation cap.\cite{series:obsatten,series:fallbacktax,series:tieredobs}
\item \textbf{Repeated-use profiling:} \texttt{profiling\_equalization} publishes a profiling/equalization summary that can be compared across releases.\cite{series:profiling}
\end{enumerate}

\begin{table}[ht]
\centering
\footnotesize
\setlength{\tabcolsep}{4pt}
\renewcommand{\arraystretch}{1.05}
\begin{tabularx}{\linewidth}{@{}l >{\raggedright\arraybackslash}X >{\raggedright\arraybackslash}X@{}}
\toprule
Line item & Interface ids (binding) & Budget meaning (one line) \\
\midrule
\texttt{primary\_path\_declaration} & \texttt{tw\_id}, \texttt{primary\_surface\_manifest\_id} & Guarded non-collusion checkpoint $(0,\rho)$ for the relayed primary path. \\
\texttt{fallback\_contact\_surface} & \texttt{tw\_id}, \texttt{state\_decl\_id}, \texttt{exposure\_nf\_id=enf.\allowbreak fallback.\allowbreak cct.\allowbreak v1} & Attenuated Tier~2 contact leakage $b_{\mathrm{att}}$ under an activation cap $p$. \\
\texttt{fallback\_tiered\_vector} & \texttt{tw\_id}, \texttt{exposure\_nf\_id=enf.\allowbreak fallback.\allowbreak tiervec.\allowbreak v1} & Tier-vector summary (Tier~0/1/2 components) for conservative composition. \\
\texttt{path\_selection\_indicator} & \texttt{tw\_id}, policy witness id & Publishes the fallback activation cap/estimate $p$ and how it is justified (policy vs audit). \\
\texttt{profiling\_equalization} & \texttt{tw\_id}, \texttt{exposure\_nf\_id=enf.\allowbreak profile.\allowbreak retry.\allowbreak v1} & Pairwise equalization / odds-inflation budget $(\eta,\delta)$ for repeated-use profiling control. \\
\bottomrule
\end{tabularx}
\caption{Five representative receipt line items in the worked bundle, shown at the ``what-to-read'' granularity. Full digests and replay hooks live in \texttt{example\_receipt.json}.}
\label{tab:lineitemglance}
\end{table}

\paragraph{Receipt excerpt (worked bundle).}
For compactness we show only a minimal excerpt. The shipped \texttt{example\_receipt.json} carries full digests and replay hooks; here we expand \texttt{fallback\_contact\_surface} to show the semantic contact-surface binding (see Synthesis~12).\cite{series:receiptitems}
\begin{verbatim}
{
  "claim_id": "...",
  "state_decl_id": "sha256:...",
  "tw_id": "sha256:...",
  "line_items": [
    {"name":"primary_path_declaration", ...},

    {"name":"fallback_contact_surface",
     "contact_surface_id":"sha256:...",
     "contact_surface":{
        "method":"GET",
        "url":"https://cid.contact/routing/v1/providers",
        "filter_protocols":["transport-bitswap"],
        "timeouts":{"http_router":"30s","routing":"60s"}
     },
     ...},

    {"name":"fallback_tiered_vector", ...},
    {"name":"path_selection_indicator", ...},
    {"name":"profiling_equalization", ...}
  ]
}
\end{verbatim}

\paragraph{What this bundle intentionally omits.}
This worked example is primarily a \emph{receipt interlock} for fallback exposure and repeated-use profiling.
A full deployment that uses delegated routing (Routing~V1, AutoConf, router overrides, caching, trustless HTTP retrieval) should also pin those control-plane knobs in \texttt{state\_decl\_id} per Synthesis~30; we keep that as an explicit extension point rather than duplicating Synthesis~30 inside this paper.\cite{series:deployedsurfaces}

\section{State-surface micro-slice}
State makes privacy claims fragile: caches, routing tables, congestion scores, and policy flags can condition the transcript and can themselves become witnesses.\cite{series:state1}
The worked example therefore publishes a \emph{minimal, diffable} state declaration (digest-bound \texttt{state\_decl\_id}) describing only what is needed to interpret the line items:
\begin{itemize}[leftmargin=*]
\item which caches and retry policies are in force,
\item which public congestion/load beacons are consulted by routing,
\item which versioned selection policy (menu) is active.
\end{itemize}
Everything else is treated as private artifact state that must be replayable by an auditor if it affects any line item.\cite{series:planstateequiv,series:auditorreplay}

\section{How the interlock is supposed to work}
This example is intentionally ``boring'': it is a mechanical contract.
\begin{enumerate}[leftmargin=*]
\item \textbf{User discovery.} A user finds the claim via the UVI/minimal interface object and obtains the receipt digest.\cite{series:userminiface}
\item \textbf{User diffing.} The user compares \texttt{receipt\_id}, \texttt{tw\_id}, \texttt{enf\_registry\_id}, and \texttt{state\_decl\_id} across releases; any semantic change forces a digest change.\cite{series:receiptitems}
\item \textbf{Auditor replay.} If challenged, the deployer provides the private artifact bundle; the auditor runs the replay plans to recompute each budget number and verify it matches the receipt.\cite{series:auditorreplay}
\item \textbf{Downstream consumption.} Endpoint tools consume the published per-lookup MaxL/odds parameters and (separately) the TW declaration for horizon reasoning.\cite{series:endpointbridge,series:accountingrulebook}
\end{enumerate}


\paragraph{Algorithmic skeleton (client-side, one lookup).}
This is intentionally high-level; concrete control-plane pinning (router set precedence, path scope, filters, caching, retrieval modes) is owned by Synthesis~30.\cite{series:deployedsurfaces}
\begin{enumerate}[leftmargin=*]
\item \textbf{Load declared state:} resolve \texttt{state\_decl\_id} to an effective router set and request-shaping defaults (including any protocol filters).
\item \textbf{Bind service interfaces:} when the receipt names a service contact surface, assemble the declared \texttt{contact\_surface} (method + URL + shaping defaults + time budgets) and verify \texttt{contact\_surface\_id} matches \texttt{sha256(JCS(contact\_surface))}.\cite{rfc8785}
\item \textbf{Attempt primary path:} query the primary service using the declared threat window and ENF; if the separation guard fails or the primary times out, record a primary failure.
\item \textbf{Gate fallback:} if fallback is taken, enforce the published activation cap (path-selection indicator) and emit only the attenuated Tier~2 vector declared in the receipt.
\item \textbf{Emit replay hooks:} store digests for the replay plan and any private artifacts (so an auditor can request the bundle and rerun the accountant).
\item \textbf{Publish receipt:} commit the tuple(s) in the release receipt so users can diff across versions.\end{enumerate}
\section{What changes across a release (micro-diffs)}

This worked bundle is meant to make \emph{semantic drift mechanically visible}.
Three common drift patterns illustrate what should (and should not) change in a receipt:

\begin{itemize}[leftmargin=*]
\item \textbf{Knob drift (same interface).} If the transcript projection is unchanged but a retry/timeout/width knob changes, the line item keeps the same \texttt{exposure\_nf\_id} but changes \texttt{knobs} and (usually) \texttt{budget\_value}. Tag \texttt{lineage=same-enf} and require replay/evidence for the new number.\cite{series:receiptitems,series:planstateequiv}
\item \textbf{State-surface drift (same mechanism, different reader set).} If a deployment silently changes which delegated-routing endpoints are contacted (AutoConf defaults, cache layers, gateway exposure), the \texttt{state\_decl\_id} must change even if the budget arithmetic is identical. Otherwise the claim is underspecified.\cite{series:deployedsurfaces,series:planstateequiv}
\item \textbf{Interface drift (new projection).} If the observable transcript grows (e.g., a new \texttt{/routing/v1} endpoint is used, or response streaming changes message timing), treat it as a new \texttt{exposure\_nf\_id} (or an explicit coarsening if the projection is strictly reduced). Do not pretend the previous certificate applies.\cite{series:traceifaces,series:receiptitems}
\end{itemize}

The equivalence/coarsening taxonomy for plan/state ids in Synthesis~18 is the intended shared language for these cases; this section is only the ``worked example'' view.\cite{series:planstateequiv}

\paragraph{User micro-diff algorithm (receipt-level).}
A user does not need to interpret every field; a deterministic diff is enough:
(1) fetch both receipts and verify their advertised digests (\texttt{receipt\_id});
(2) compare the small set of top-level interface/state ids;
(3) compare line items by \texttt{name}.
\begin{verbatim}
Top-level: tw_id, enf_registry_id, state_decl_id,
           primary_surface_manifest_id, replay_plans_id
Per line item (by name): exposure_nf_id, knobs, budget_value,
                         evidence_id, replay_hook, lineage
\end{verbatim}
Interpretation: changing \texttt{exposure\_nf\_id} is interface drift; changing \texttt{state\_decl\_id} is control-plane/reader-set drift; changing \texttt{knobs} or \texttt{budget\_value} with the same \texttt{exposure\_nf\_id} is knob drift (same interface, new certificate).

\paragraph{Canonicalization notes.}
Digest IDs should be computed over canonical bytes, not ad hoc pretty-printed JSON.
For JSON objects we recommend the JSON Canonicalization Scheme (JCS, RFC~8785), i.e., define \texttt{sha256(JCS(obj))}.\cite{rfc8785}
For human micro-diffs between two receipts-as-JSON, also canonicalize \emph{semantic sets} (sort and deduplicate arrays that represent sets, such as protocol filter lists or router allowlists) and normalize duration strings (parse \texttt{30s}/\texttt{0.5m} into a single unit) so reordering/formatting noise does not masquerade as drift.

\paragraph{Canned drift cases (bundle-provided).}
The shipped bundle includes \texttt{artifacts/example\_drift\_cases.json}, a three-row catalog of ``what changed'' successors for the base release label \texttt{worked-example-draft-59}. It is intended as a sanity check for the taxonomy above (refresh-only vs replay-required vs recertification).
\begin{table}[t]
\centering
\begin{tabularx}{\linewidth}{@{}l l X@{}}
\toprule
Case & Class & What changed / required action \\
\midrule
\texttt{refresh-only-A} & \texttt{refresh\_only} & Log checkpoint / user-watch text only; rebind receipt lineage, no budget re-evaluation. \\
\texttt{replay-required-B} & \texttt{replay\_changed\_surfaces} & Fallback observation upper bound tightened ($p_{\mathrm{obs}}:0.02\to0.015$); rerun declared evaluators and publish a fresh notarized comparison. \\
\texttt{full-recertification-C} & \texttt{full\_recertification} & Threat window and state declaration (cache trigger) changed; treat as certified-interface drift and require recertification gating. \\
\bottomrule
\end{tabularx}
\caption{Three canned drift cases shipped in the worked-example bundle.}\label{tab:workedexample-driftcases}
\end{table}

\paragraph{Example output (from the shipped compare report).}
The bundle also includes \texttt{artifacts/example\_compare\_report.json}; its \texttt{user\_fast\_diff} contains a ``first alarm'' message and the minimal numeric fields that casual users usually care about. For the canned successor in that report, the fallback budget improves while \texttt{exposure\_nf\_id} stays fixed:
\begin{verbatim}
base:     budget_value=0.05635, total_summary_bits=0.26590
successor budget_value=0.04240, total_summary_bits=0.26518
(same exposure_nf_id; receipt digest changed)
\end{verbatim}

In this example, \emph{all} control-plane drift that changes the effective reader set or request shaping (endpoint sets, AutoConf expansion rules, environment overrides, protocol filters, retrieval-mode switches) is treated as \textbf{state-surface drift} and therefore must bump \texttt{state\_decl\_id} even if the budget arithmetic is unchanged.\cite{series:deployedsurfaces,series:planstateequiv}

\paragraph{Three tiny examples.}
\begin{enumerate}[leftmargin=*]
\item \textbf{Window drift:} keep label \texttt{tw.reader\_24h.q500.v1} but change $Q$ (or the reset rule). The digest $TW$ changes and downstream $Q$-scale reasoning must be revisited.\cite{series:threatwindows}
\item \textbf{Surface drift:} keep label \texttt{enf.fallback.cct.v1} but add a new Tier~2 field. The digest $E_i$ changes; a user diff catches that the exposure surface widened.\cite{series:traceifaces}
\item \textbf{Evidence drift:} keep the profiling label but change the audit projection (bucket widths) or the confidence knob. The \texttt{evidence\_id} digest changes and auditors replay via the \texttt{replay\_hook}.\cite{series:receiptitems,series:auditorreplay}
\end{enumerate}




\paragraph{Request-shaping drift (filters and budgets).}
If the effective delegated-routing request adds, removes, or changes request-shaping defaults or budgets (protocol filters, router allow/deny policy, path scope, routing/retrieval timeouts that govern downgrade and fallback), treat it as \textbf{state-surface drift}: it changes \emph{who is contacted} and what termination-time behavior is feasible, and therefore can change fallback rate and exposure.
Bump \texttt{state\_decl\_id} and treat it as recertification-triggering.\cite{series:deployedsurfaces,series:planstateequiv}

\section{Illustrative numerical slice}
\label{sec:numerical}
This note includes exactly one numerical ``shape'' so a reader sees how a receipt number relates to a projection and an activation cap.
All arithmetic is recomputed by the shipped validator for the maintained bundle cut (\texttt{example\_validation\_report.json}).\cite{series:auditorreplay}

\paragraph{Attenuation by activation probability.}
When a Tier~2 projection leaks $b$ bits \emph{conditional on fallback occurring}, and fallback occurs with probability at most $p$ within the declared threat window, the published attenuated leakage is
\[
b_{\mathrm{att}} \;=\; \log_2\!\big((1-p) + p\,2^b\big),
\]
the standard ``odds inflation'' attenuation used throughout the series.\cite{series:oddsinflation,series:obsatten}
The worked bundle instantiates this for the fallback contact surface line item; the receipt tags whether $p$ is a policy cap or an empirical estimate.

\begin{center}
\small
\begin{tabular}{@{}r r r@{}}
\toprule
$b$ (bits, conditional) & $p$ (fallback cap) & $b_{\mathrm{att}}=\log_2((1-p)+p2^b)$ (bits)\\
\midrule
8 & 0.05 & 3.78\\
8 & 0.01 & 1.83\\
12 & 0.01 & 5.39\\
\bottomrule
\end{tabular}
\end{center}
\vspace{-0.5em}

\noindent These rows are \emph{only} intuition for how the attenuation behaves; the receipt’s authoritative numbers are rederived by the validator from the shipped bundle objects.

\paragraph{Why this is enough for a worked example.}
The details---how $p$ is obtained (selection taxes), how $b$ is computed for a concrete ENF (trace interfaces), and how repeated audits spend $\delta_{\mathrm{stat}}$---are each owned by separate synthesis notes.\cite{series:fallbacktax,series:traceifaces,series:stataudit,series:anytimespend}
Keeping this section minimal avoids re-deriving results that already have their own ``referee-grade'' homes.

\section{Optional variants (two one-paragraph receipts)}
\paragraph{Prefix-fetch equalization projection.}
A variant receipt replaces the fallback ENF with a prefix-fetch projection (coarser query privacy via requesting a bucket/prefix rather than a single key).
The interface change is surfaced \emph{only} by a different ENF digest and different evidence digest; the rest of the receipt machinery is unchanged.\cite{series:prefixcapacities,series:stataudit}

\paragraph{Coarsening a large-alphabet projection.}
A second variant receipt demonstrates safe coarsening of a large-alphabet contact projection into an audit-feasible partition.
Again, the only user-visible change is via the ENF and evidence digests, plus the resulting budget number.\cite{series:coarsening}
If the partition map is selected from data (rather than fixed ex ante), the selection discipline note applies: treat the map family as witness-indexed or separate pilot selection from frozen certification under change control.\cite{series:selectiondiscipline}

\section{Takeaways and open edges (short)}
\begin{itemize}[leftmargin=*]
\item A referee-facing object can stay tiny: digests + a handful of line items + replay hooks.
\item ``Same name, different semantics'' is handled by digest IDs (ENF/TW/state), not by prose.
\item Notation (keys/state/transcripts, projections) is standardized in Synthesis~8; this note only instantiates it.\cite{series:notation}
\item Guarded endpoints are essential: a primary-path claim that depends on non-collusion must ship its $\rho$.
\item Tier vectors and selection tax compile to a single per-lookup MaxL (or odds) parameter that downstream tools can consume.\cite{series:endpointbridge}
\item Remaining research: robust horizon accounting under stateful dependence and operator-safe drift discipline; the state series and accounting rulebook are the intended import points.\cite{series:state1,series:accountingrulebook}
\end{itemize}

\begin{thebibliography}{99}\small


\bibitem{rfc8785}
A.~Rundgren.
\newblock JSON Canonicalization Scheme (JCS).
\newblock RFC~8785, IETF, 2020.

\bibitem{series:artifactpolicy}
Anonymous.

\newblock Artifact and Disclosure Policy for the One-True-Archive.

\newblock Synthesis~6, The-One-True-Archive (2026).

\bibitem{series:deployedsurfaces}
Anonymous.

\newblock Deployed Observation Surfaces for Anonymous DHTs: Control-plane knobs in IPFS/libp2p delegated routing and OHTTP evolution.

\newblock Series synthesis note (Synthesis~30), 2026. (This archive.)

\bibitem{series:receiptitems}
Anonymous.

\newblock Receipt Line Items for Anonymous-DHT Privacy Claims: Minimal Tuple Schema and Drift Semantics.

\newblock Series synthesis note (Synthesis~12), 2026. (This archive.)

\bibitem{series:notation}
Anonymous.

\newblock Notation and a Minimal Model for the Series.

\newblock Series synthesis note (Synthesis~8), 2026. (This archive.)

\bibitem{series:traceifaces}
Anonymous.

\newblock Trace Interfaces for Anonymous DHTs.

\newblock Synthesis~3, The-One-True-Archive (2026).

\bibitem{series:threatwindows}
Anonymous.

\newblock Threat Models and Repetition Windows for Anonymous DHTs.

\newblock Synthesis~4, The-One-True-Archive (2026).

\bibitem{series:endpointbridge}
Anonymous.

\newblock Endpoint Metrics Bridge: Odds Inflation, PML/MaxL, and Identification Sizing.

\newblock Synthesis~5, The-One-True-Archive (2026).

\bibitem{series:oddsinflation}
Anonymous.

\newblock Odds Inflation, Privacy Loss, and Endpoint Anonymity Bounds.

\newblock Anonymity~A, The-One-True-Archive (2026).

\bibitem{series:profiling}
Anonymous.

\newblock Anonymous DHT Profiling and the Equalization Contract.

\newblock Anonymity~B, The-One-True-Archive (2026).

\bibitem{series:primarysepmanifest}
Anonymous.

\newblock Primary-Path Separation Manifests: Guard Objects for Non-Collusion Assumptions in Reader-Privacy Receipts.

\newblock Synthesis~22, The-One-True-Archive (2026).

\bibitem{series:probsep}
Anonymous.

\newblock Probabilistic Separation Budgets: Separation Failure as a Receipt-Grade $(\varepsilon,\delta)$ Adjunct.

\newblock Synthesis~23, The-One-True-Archive (2026).

\bibitem{series:obsatten}
Anonymous.

\newblock Observation Attenuation for Repeated-Use Privacy Budgets.

\newblock Synthesis~11, The-One-True-Archive (2026).


\bibitem{series:planstateequiv}
Anonymous.

\newblock Digest-Bound Replay Plans and State Declarations: Equivalence, Minimality, and Change Control.

\newblock Synthesis~18, The-One-True-Archive (2026).

\bibitem{series:userminiface}
Anonymous.

\newblock User-Verifiable Receipt Interfaces.

\newblock Synthesis~13, The-One-True-Archive (2026).

\bibitem{series:fallbacktax}
Anonymous.

\newblock Adaptive Fallback as a Witness Stream: The Path-Selection Tax.

\newblock Synthesis~14, The-One-True-Archive (2026).

\bibitem{series:tieredobs}
Anonymous.

\newblock Tiered Observation Vectors for Anonymous-DHT Receipts.

\newblock Synthesis~15, The-One-True-Archive (2026).

\bibitem{series:abomoinl}
Anonymous.

\newblock ABOM and OINL for Anonymity Claims.

\newblock Anonymity series Paper~C, The-One-True-Archive (2026).

\bibitem{series:bossfightB}
Anonymous.

\newblock AnonDHT in the Boss Fight Regime: Compiling a MaxL Budget into Lookup Knobs.

\newblock Boss Fight series Paper~B, The-One-True-Archive (2026).


\bibitem{series:prefixcapacities}
Anonymous.

\newblock Prefix Capacities and Separation Cuts for Shared-Kernel Termination-Time Hiding.

\newblock Published mathematics note, 2026-01-25. (This archive: \texttt{published/2026-01-25\_prefix\_capacities}.)

\bibitem{series:stoptimenf}
Anonymous.

\newblock Stop-Time Padding Addendum, Max-Mixture Separations and Tail-Sign Deadline Normal Forms.

\newblock Published mathematics addendum, 2026-01-26. (This archive: \texttt{published/2026-01-26\_stop\_time\_padding\_addendum}.)

\bibitem{series:state1}
Anonymous.

\newblock State-Dependent Anonymity for Anonymous DHTs.

\newblock AnonDHT state series Paper~1, The-One-True-Archive (2026).

\bibitem{series:condcert}
Anonymous.

\newblock Conditional Per-Step Budget Certificates: A Receipt-Grade Interface for Adaptive Horizon Claims.

\newblock Synthesis~21, The-One-True-Archive (2026).

\bibitem{series:accountingrulebook}
Anonymous.

\newblock Receipt Accounting and Horizon Composition for Anonymous-DHT Budgets.

\newblock Synthesis~19, The-One-True-Archive (2026).


\bibitem{series:stataudit}
Anonymous.

\newblock Discrete Equalization Audits from Logs: Menu-Level Finite-Sample TV Certificates.

\newblock Synthesis~25, The-One-True-Archive (2026).

\bibitem{series:anytimespend}
Anonymous.

\newblock Anytime Statistical Spending for Menu Audits: Change-Control-Safe Confidence Budgets.

\newblock Synthesis~26, The-One-True-Archive (2026).

\bibitem{series:coarsening}
Anonymous.

\newblock Coarsening and Large-Alphabet Projections for Log Audits: Safe Partitioning Rules that Keep Certificates Sample-Feasible.

\newblock Synthesis~28, The-One-True-Archive (2026).

\bibitem{series:selectiondiscipline}
Anonymous.

\newblock Selection Discipline for Data-Dependent Bucket Maps.

\newblock Synthesis~29, The-One-True-Archive (2026).

\bibitem{series:auditorreplay}
Anonymous.

\newblock Auditor Replay for Anonymous-DHT Receipt Line Items: A Mechanical Checklist for Digest-Bound Claims.

\newblock Series synthesis note (Synthesis~20), 2026. (This archive.)



\end{thebibliography}

\end{document}
